%%% FIELD proof-uniform-minimax-frontier
Put
\[
L_{d,\alpha}=\log\frac{d+1}{\alpha},\qquad
K=\left\lceil 8\log\frac{2d}{\alpha}\right\rceil,
\qquad q=\left\lfloor\frac nK\right\rfloor,
\qquad \rho=\sqrt{\frac{8V_\psi}{q}},
\]
as in Definition \(\mathrm{def:chart\text{-}moment\text{-}rule}\).  Fix
\((P,\theta)\in\mathcal M_d^{\mathrm H}(\lambda)\), and define
\[
\mu_j=\mathbb E_P\psi(X_j,Y_j).
\]
Lemma \(\mathrm{lem:subhalf\text{-}chart\text{-}calibration}\) gives, for every \(j\in I\),
\[
\mu_j=\theta_j\,2t_j(1-2a_j),
\qquad
\operatorname{Var}_P\{\psi(X_j,Y_j)\}\le V_\psi,
\qquad
H_{j,-\theta_j}\le C_Ht_j^2.
\tag{1}
\]
In particular, Definition \(\mathrm{def:subhalf\text{-}hermite\text{-}class}\) imposes \(t_j>0\) and \(a_j\le a_+<1/2\), so
\[
\operatorname{sign}(\mu_j)=\theta_j.
\tag{2}
\]

Suppose first that \(n\ge2K\).  For a block \(A_k\), let
\[
\overline\psi_{jk}=q^{-1}\sum_{i\in A_k}\psi_{ij}.
\]
The independent-sample model and (1) imply that this block mean has mean \(\mu_j\) and variance at most \(V_\psi/q\).  Chebyshev's inequality, obtained directly by applying Markov's inequality to \((\overline\psi_{jk}-\mu_j)^2\), gives
\[
P^{\otimes n}\bigl(|\overline\psi_{jk}-\mu_j|>\rho\bigr)
\le \frac{V_\psi}{q\rho^2}=\frac18.
\tag{3}
\]
The \(K\) block means are independent because the observation-index blocks are disjoint.  If their lower median differs from \(\mu_j\) by more than \(\rho\), at least \(K/2\) block means satisfy the event in (3).  A union bound over subsets of blocks therefore yields
\[
\begin{aligned}
P^{\otimes n}\bigl(|\widehat\mu_j-\mu_j|>\rho\bigr)
&\le \sum_{s=\lceil K/2\rceil}^{K}\binom Ks(1/8)^s\\
&\le 2^K(1/8)^{K/2}
=2^{-K/2}\\
&\le \exp\left\{-4(\log2)\log\frac{2d}{\alpha}\right\}
\le \frac{\alpha}{2d}.
\end{aligned}
\tag{4}
\]
For the last inequality, \(K\ge8\log(2d/\alpha)\), \(4\log2>1\), and \(\alpha/(2d)<1\).  Thus a union bound over the deterministic set \(I\) gives
\[
P^{\otimes n}(\mathcal E)\ge1-\alpha/2,
\qquad
\mathcal E=\left\{\max_{j\in I}|\widehat\mu_j-\mu_j|\le\rho\right\}.
\tag{5}
\]
This probability bound is uniform over the chart class because the variance bound in (1) is uniform.

Membership in \(\mathcal M_d^{\mathrm H}(\lambda)\) entails membership in \(\mathcal M_d(\lambda)\), hence \(H_{j,-\theta_j}\ge\lambda\).  Combining this fact with (1) gives
\[
t_j^2\ge\frac{\lambda}{C_H},
\qquad
|\mu_j|^2=4t_j^2(1-2a_j)^2
\ge \frac{4(1-2a_+)^2}{C_H}\lambda.
\tag{6}
\]
When \(n\ge2K\), \(q=\lfloor n/K\rfloor\ge n/(2K)\), and therefore
\[
\rho^2\le\frac{16V_\psi K}{n}.
\tag{7}
\]
For \(d\ge1\) and \(0<\alpha\le1/4\),
\[
\log\frac{2d}{\alpha}\le2L_{d,\alpha},
\qquad
K\le9\log\frac{2d}{\alpha}\le18L_{d,\alpha}.
\tag{8}
\]
Choose \(C\), depending only on \(V_\psi\), \(C_H\), and the fixed endpoint \(a_+\), so that
\[
C>\max\left\{36,\frac{288V_\psi C_H}{(1-2a_+)^2}\right\}.
\tag{9}
\]
Because \(0<\lambda\le\lambda_H\le1\), the inequality \(n\lambda\ge CL_{d,\alpha}\), together with (8)--(9), implies \(n\ge2K\).  Equations (6)--(9) also imply
\[
4\rho^2<\min_{j\in I}|\mu_j|^2.
\tag{10}
\]
Indeed, (7)--(8) give \(4\rho^2\le1152V_\psi L_{d,\alpha}/n\le1152V_\psi\lambda/C\), while (6) bounds the right-hand side of (10) below by \(4(1-2a_+)^2\lambda/C_H\), and (9) compares these constants.

On \(\mathcal E\), (10) implies that each interval \([\widehat\mu_j-\rho,\widehat\mu_j+\rho]\) lies strictly in the open half-line with sign \(\theta_j\).  Definition \(\mathrm{def:chart\text{-}moment\text{-}rule}\) then gives \(B_j=\{\theta_j\}\) for every \(j\), and consequently
\[
\widehat{\mathcal C}_n^{\mathrm H}=\{\theta\},
\qquad
\widehat\theta_n^{\mathrm H}=\theta
\quad\text{on }\mathcal E.
\]
Together with (5), this proves
\[
\sup_{(P,\theta)\in\mathcal M_d^{\mathrm H}(\lambda)}
P^{\otimes n}(\widehat\theta_n^{\mathrm H}\ne\theta)
\le\alpha/2.
\tag{11}
\]

Lemma \(\mathrm{lem:chart\text{-}product\text{-}testing\text{-}lower}\) supplies \(c_L>0\) such that every measurable estimator has worst-case error greater than \(\alpha\) whenever \(n\lambda\le c_LL_{d,\alpha}\).  Combining that lower bound with (11) and Definition \(\mathrm{def:chart\text{-}minimax\text{-}frontier}\), and absorbing the integer ceiling into the constants using \(L_{d,\alpha}\ge\log8\) and \(\lambda\le1\), yields constants \(0<c<C<\infty\) for which
\[
c\frac{L_{d,\alpha}}{\lambda}
\le N_{\mathrm H}^*(d,\alpha,\lambda)
\le C\frac{L_{d,\alpha}}{\lambda}.
\]
Finally, Definition \(\mathrm{def:subhalf\text{-}hermite\text{-}class}\) gives \(\mathcal M_d^{\mathrm H}(\lambda)\subseteq\mathcal M_d(\lambda)\).  The supremum risk over \(\mathcal M_d(\lambda)\) is therefore no smaller than the supremum over its chart subclass, so the same lower bound holds for \(N^*(d,\alpha,\lambda)\); this also agrees with Theorem \(\mathrm{thm:product\text{-}testing\text{-}converse}\).  Lemma \(\mathrm{lem:peters\text{-}disjoint\text{-}pair\text{-}specialization}\), applied to the member conditions in Definition \(\mathrm{def:base\text{-}model}\), ensures that \(\theta\) is the observationally identified causal orientation on the disjoint-edge restricted-ANM domain.  Neither the chart estimator nor its upper bound is transferred to the larger class.
%%% FIELD proof-simultaneous-graph-confidence-set
All block boundaries in Definition \(\mathrm{def:chart\text{-}moment\text{-}rule}\) are deterministic.  Each block average is a Borel function of the sample, a finite lower order statistic is Borel measurable, and every condition defining \(B_j\) is a Borel inequality.  Since \(\{-1,+1\}^I\) is finite, the Cartesian product \(\widehat{\mathcal C}_n^{\mathrm H}=\prod_{j\in I}B_j\) is a measurable random set.  The rule evaluates the fixed observable functional \(\psi\) once for each pair coordinate and observation, hence it uses exactly \(d\) coordinatewise moment queries.  Its \(nd\) evaluations and block sums require \(O(dn)\) arithmetic operations.  A deterministic linear-time selection algorithm computes each median in \(O(K)\) operations, and \(K\le n/2\) because \(n\ge2K\), so the total cost remains \(O(dn)\).

Fix \((P,\theta)\in\mathcal M_d^{\mathrm H}(\lambda)\) and write
\[
\mu_j=\mathbb E_P\psi(X_j,Y_j).
\]
The block argument in the proof of Theorem \(\mathrm{thm:uniform\text{-}minimax\text{-}frontier}\), using the uniform variance bound in Lemma \(\mathrm{lem:subhalf\text{-}chart\text{-}calibration}\), proves
\[
P^{\otimes n}\left\{\max_{j\in I}|\widehat\mu_j-\mu_j|
\le\rho_{n,d,\alpha}\right\}
\ge1-\alpha/2.
\tag{12}
\]
On the event in (12), the interval \([\widehat\mu_j-\rho_{n,d,\alpha},\widehat\mu_j+\rho_{n,d,\alpha}]\) contains \(\mu_j\).  The calibration lemma gives
\[
\mu_j=\theta_j\,2t_j(1-2a_j).
\]
Definition \(\mathrm{def:subhalf\text{-}hermite\text{-}class}\) gives \(t_j>0\) and \(a_j\le a_+<1/2\), so \(\operatorname{sign}(\mu_j)=\theta_j\).  If \(\theta_j=+1\), interval containment gives \(\widehat\mu_j+\rho_{n,d,\alpha}\ge\mu_j>0\), and hence \(+1\in B_j\).  If \(\theta_j=-1\), it gives \(\widehat\mu_j-\rho_{n,d,\alpha}\le\mu_j<0\), and hence \(-1\in B_j\).  Thus \(\theta_j\in B_j\) in both cases, and (12) proves
\[
\inf_{(P,\theta)\in\mathcal M_d^{\mathrm H}(\lambda)}
P^{\otimes n}\left\{\theta\in\widehat{\mathcal C}_n^{\mathrm H}\right\}
\ge1-\alpha/2.
\tag{13}
\]
Because \(t_j>0\), the chart mechanism is non-affine.  Lemma \(\mathrm{lem:population\text{-}direction\text{-}margin}\) therefore gives \(H_{j,\theta_j}=0<H_{j,-\theta_j}\).  Hence the sign used above recovers the causal orientation identified by the population residual-HSIC functional; it does not define a different target.

For the constant \(C\) in Theorem \(\mathrm{thm:uniform\text{-}minimax\text{-}frontier}\), the implication \(n\lambda\ge C\log((d+1)/\alpha)\) gives
\[
2\rho_{n,d,\alpha}<\min_{j\in I}|\mu_j|.
\tag{14}
\]
On the event in (12), (14) places every interval strictly in the half-line with sign \(\theta_j\).  Thus \(B_j=\{\theta_j\}\) for every \(j\), and
\[
\widehat{\mathcal C}_n^{\mathrm H}=\{\theta\}.
\]
Combining this identity with (12) proves the asserted uniform singleton probability.  Finally, Lemma \(\mathrm{lem:chart\text{-}product\text{-}testing\text{-}lower}\) states, on the same chart class and below its constant-factor threshold, that no random set can have both uniform coverage and uniform singleton probability at least \(1-\alpha/2\).  This proves the converse assertion.
%%% FIELD proof-single-pair-reduction
When \(d=1\), one has \(I=\{1\}\), so the Cartesian product in Definition \(\mathrm{def:chart\text{-}moment\text{-}rule}\) is
\[
\widehat{\mathcal C}_n^{\mathrm H}=\prod_{j\in I}B_j=B_1.
\]
Substituting \(j=1\) into that definition gives
\[
+1\in B_1
\quad\Longleftrightarrow\quad
\widehat\mu_1+\rho_{n,1,\alpha}\ge0,
\qquad
-1\in B_1
\quad\Longleftrightarrow\quad
\widehat\mu_1-\rho_{n,1,\alpha}\le0,
\]
which is the asserted pairwise inversion.

For every member of \(\mathcal M_1^{\mathrm H}(\lambda)\), Lemma \(\mathrm{lem:subhalf\text{-}chart\text{-}calibration}\) gives
\[
\mathbb E_P\psi(X_1,Y_1)=\theta_1\,2t_1(1-2a_1).
\]
The sub-half restrictions impose \(t_1>0\) and \(a_1\le a_+<1/2\), so the sign of this expectation is \(\theta_1\).  Since \(t_1>0\), the chart mechanism is non-affine.  Lemma \(\mathrm{lem:population\text{-}direction\text{-}margin}\) then yields
\[
H_{1,\theta_1}=0<H_{1,-\theta_1},
\]
showing that the moment sign equals the unique residual-HSIC-identified causal direction.

Finally, specialize the two-sided conclusion of Theorem \(\mathrm{thm:uniform\text{-}minimax\text{-}frontier}\) to \(d=1\).  Since \(\log((d+1)/\alpha)=\log(2/\alpha)\), constants independent of \(\alpha\) and \(\lambda\) satisfy
\[
c\lambda^{-1}\log\frac2\alpha
\le N_{\mathrm H}^*(1,\alpha,\lambda)
\le C\lambda^{-1}\log\frac2\alpha,
\qquad
0<\alpha\le1/4,\quad 0<\lambda\le\lambda_H.
\]
This is the asserted \(\asymp\) relation.  The coverage interpretation of the displayed inversion follows by specializing Theorem \(\mathrm{thm:simultaneous\text{-}graph\text{-}confidence\text{-}set}\) to \(d=1\).
